\chapter{Notation}
\section{Special operators and symbols}
\subsection*{Matrix diagonal extraction}
Operator $\diag$ extracts the elements on the
main diagonal of a matrix, i.e.\
$\diag\T{A} = a_{ii}$.

\subsection*{Matrix direct sum}
A direct sum of two matrices is defined as
\begin{equation}
  \T{A} \oplus \T{B} =
  \begin{bmatrix}
    \T{A} & \T{0} \\
    \T{0} & \T{B}
  \end{bmatrix}.
\end{equation}
Furthermore, a sum of $n$ elements is defined as
\begin{equation}
  \bigoplus_{i = 1}^n \T{A}_i =
  \begin{bmatrix}
    \T{A}_1 & \T{0}   & \cdots & \T{0}   \\
    \T{0}   & \T{A}_2 & \cdots & \T{0}   \\
    \vdots  & \vdots  & \ddots & \vdots  \\
    \T{0}   & \T{0}   & \cdots & \T{A}_n
  \end{bmatrix}.
\end{equation}

\subsection*{Inner product}
A bilinear form $\iprd[\T{A}]{\,\cdot\,}{\,\cdot\,}  : \Complex^n \times \Complex^n \to \Complex$
named $\T{A}$-inner product is defined as $\iprd[\T{A}]{\T{u}}{\T{v}} = \T{u}^* \T{A} \T{v}$.
In the special case of $\T{A} = \T{I}$, which is denoted by simply
$\iprd{\,\cdot\,}{\,\cdot\,}$, the
inner product is equivalent to the dot product of two vectors:
$\iprd[\T{I}]{\T{u}}{\T{v}} = \iprd{\T{u}}{\T{v}} = \T{u}^* \T{v}$.

\section{\nameCaSymbols}

\subsubsection*{Typographical symbols}
\S\ chapter, section, subsection

% $\blacksquare$ end of a proof
% 
% $\vartriangle$ end of a remark or example
% 
% $\blacktriangle$ end of a definition

\subsubsection*{Mathematical symbols}

$j$ imaginary unit

$\T{I}$ identity matrix

$\T{0}$ zero matrix

$\T{P}$ permutation matrix

$\T{Q}$, $\T{U}$ unitary (orthogonal) matrix

$\T{T}$ transformation matrix

$\T{\mathit{R}}(\phi)$ rotation matrix

$\mathbb{R}^n$ n-dimensional real space

$\SO{3}$ special orthogonal group of three-dimensional rotations

$\so{3}$ Lie algebra of $\SO{3}$

$\SE{3}$ special Euclidean group of three-dimensional rigid-body motions

$\se{3}$ Lie algebra of $\SE{3}$

${}^{A}\T{R}_{B}$ rotation matrix that maps coordinates from frame $B$ to frame $A$

${}^{A}\T{p}_{B}$ position of the origin of frame $B$ expressed in frame $A$

${}^{A}\T{T}_{B}$ homogeneous transformation from frame $B$ to frame $A$

\begin{equation}
  {}^{A}\T{T}_{B} =
  \begin{bmatrix}
    {}^{A}\T{R}_{B} & {}^{A}\T{p}_{B} \\
    \T{0}_{1\times 3} & 1
  \end{bmatrix}.
\end{equation}

$\T{q} = [q_w, q_x, q_y, q_z]^{\mathsf{T}}$ unit quaternion, with
$q_w^2 + q_x^2 + q_y^2 + q_z^2 = 1$

$\overline{\T{q}}$ quaternion conjugate

$\hatmap{\T{v}}$ skew-symmetric matrix associated with vector $\T{v}$

\subsubsection*{Craig modified Denavit--Hartenberg parameters}

$a_{i-1}$ link length

$\alpha_{i-1}$ link twist

$d_i$ link offset

$\theta_i$ joint angle

In the Craig modified Denavit--Hartenberg convention, the transform from frame
$i$ to frame $i-1$ is
\begin{equation}
  {}^{i-1}\T{T}_{i} =
  \mathrm{Rot}_{x}(\alpha_{i-1})\,
  \mathrm{Trans}_{x}(a_{i-1})\,
  \mathrm{Rot}_{z}(\theta_i)\,
  \mathrm{Trans}_{z}(d_i).
\end{equation}


\section{\nameCaAcronyms}
\label{sc:abbreviationsAndAcronyms}
%\vspace{-1 cm}
%\printacronyms[name = {}]
\printacronyms[name={}, heading=none]
